\documentclass[10pt,a4paper]{amsart}
\usepackage{amsmath,amssymb,amsthm,graphicx}
% For commutative diagrams we use Xy-pic
\usepackage[all]{xy}

% We use 2cell for 2-commutative diagrams.
\xyoption{2cell}
% Theorem environments.
%
\theoremstyle{plain}
\newtheorem{theorem}[subsection]{Theorem}
\newtheorem{proposition}[subsection]{Proposition}
\newtheorem{lemma}[subsection]{Lemma}

\theoremstyle{definition}
\newtheorem{definition}[subsection]{Definition}
\newtheorem{example}[subsection]{Example}
\newtheorem{exercise}[subsection]{Exercise}
\newtheorem{situation}[subsection]{Situation}

\theoremstyle{remark}
\newtheorem{remark}[subsection]{Remark}
\newtheorem{remarks}[subsection]{Remarks}

\numberwithin{equation}{subsection}

% Macros
%
\def\lim{\mathop{\mathrm{lim}}\nolimits}
\def\colim{\mathop{\mathrm{colim}}\nolimits}
\def\Spec{\mathop{\mathrm{Spec}}}
\def\Hom{\mathop{\mathrm{Hom}}\nolimits}
\def\Ext{\mathop{\mathrm{Ext}}\nolimits}
\def\SheafHom{\mathop{\mathcal{H}\!\mathit{om}}\nolimits}
\def\SheafExt{\mathop{\mathcal{E}\!\mathit{xt}}\nolimits}
\def\Sch{\mathit{Sch}}
\def\Mor{\mathop{\mathrm{Mor}}\nolimits}
\def\Ob{\mathop{\mathrm{Ob}}\nolimits}
\def\Sh{\mathop{\mathit{Sh}}\nolimits}
\def\NL{\mathop{N\!L}\nolimits}
\def\CH{\mathop{\mathrm{CH}}\nolimits}
\def\proetale{{pro\text{-}\acute{e}tale}}
\def\etale{{\acute{e}tale}}
\def\QCoh{\mathit{QCoh}}
\def\Ker{\mathop{\mathrm{Ker}}}
\def\Im{\mathop{\mathrm{Im}}}
\def\Coker{\mathop{\mathrm{Coker}}}
\def\Coim{\mathop{\mathrm{Coim}}}

% Boxtimes
%
\DeclareMathSymbol{\boxtimes}{\mathbin}{AMSa}{"02}

%
% Macros for moduli stacks/spaces
%
\def\QCohstack{\mathcal{QC}\!\mathit{oh}}
\def\Cohstack{\mathcal{C}\!\mathit{oh}}
\def\Spacesstack{\mathcal{S}\!\mathit{paces}}
\def\Quotfunctor{\mathrm{Quot}}
\def\Hilbfunctor{\mathrm{Hilb}}
\def\Curvesstack{\mathcal{C}\!\mathit{urves}}
\def\Polarizedstack{\mathcal{P}\!\mathit{olarized}}
\def\Complexesstack{\mathcal{C}\!\mathit{omplexes}}
% \Pic is the operator that assigns to X its picard group, usage \Pic(X)
% \Picardstack_{X/B} denotes the Picard stack of X over B
% \Picardfunctor_{X/B} denotes the Picard functor of X over B
\def\Pic{\mathop{\mathrm{Pic}}\nolimits}
\def\Picardstack{\mathcal{P}\!\mathit{ic}}
\def\Picardfunctor{\mathrm{Pic}}
\def\Deformationcategory{\mathcal{D}\!\mathit{ef}}
\usepackage{cleveref}
\crefname{originalreference}{original reference}{original references}
\Crefname{originalreference}{Original reference}{Original references}
\makeatletter
\newlabel{item-surjective}{{06AP}{source}{}{sourcepointer-1}{}}
\newlabel{item-surjective@cref}{{[originalreference][0][]06AP}{[1][0][]source}{}{}{}}
\makeatother
\crefname{originalreference}{original reference}{original references}
\Crefname{originalreference}{Original reference}{Original references}
\makeatletter
\newlabel{item-vanishing}{{0A53}{source}{}{sourcepointer-2}{}}
\newlabel{item-vanishing@cref}{{[originalreference][0][]0A53}{[1][0][]source}{}{}{}}
\makeatother
\crefname{originalreference}{original reference}{original references}
\Crefname{originalreference}{Original reference}{Original references}
\makeatletter
\newlabel{lemma-compatible-shriek-push-finite}{{0GJ1}{source}{}{sourcepointer-3}{}}
\newlabel{lemma-compatible-shriek-push-finite@cref}{{[originalreference][0][]0GJ1}{[1][0][]source}{}{}{}}
\makeatother
\crefname{originalreference}{original reference}{original references}
\Crefname{originalreference}{Original reference}{Original references}
\makeatletter
\newlabel{lemma-finite-pushforward-statements}{{0A5D}{source}{}{sourcepointer-4}{}}
\newlabel{lemma-finite-pushforward-statements@cref}{{[originalreference][0][]0A5D}{[1][0][]source}{}{}{}}
\makeatother
\crefname{originalreference}{original reference}{original references}
\Crefname{originalreference}{Original reference}{Original references}
\makeatletter
\newlabel{lemma-pullback-filtered}{{0A3R}{source}{}{sourcepointer-5}{}}
\newlabel{lemma-pullback-filtered@cref}{{[originalreference][0][]0A3R}{[1][0][]source}{}{}{}}
\makeatother
\crefname{originalreference}{original reference}{original references}
\Crefname{originalreference}{Original reference}{Original references}
\makeatletter
\newlabel{lemma-ses-statements}{{0A5C}{source}{}{sourcepointer-6}{}}
\newlabel{lemma-ses-statements@cref}{{[originalreference][0][]0A5C}{[1][0][]source}{}{}{}}
\makeatother
\crefname{originalreference}{original reference}{original references}
\Crefname{originalreference}{Original reference}{Original references}
\makeatletter
\newlabel{lemma-somewhat-easier}{{0A3Q}{source}{}{sourcepointer-7}{}}
\newlabel{lemma-somewhat-easier@cref}{{[originalreference][0][]0A3Q}{[1][0][]source}{}{}{}}
\makeatother
\crefname{originallemma}{original lemma}{original lemmas}
\Crefname{originallemma}{Original lemma}{Original lemmas}
\makeatletter
\newlabel{more-morphisms-lemma-quasi-finite-separated-pass-through-finite}{{05K0}{source}{}{sourcepointer-8}{}}
\newlabel{more-morphisms-lemma-quasi-finite-separated-pass-through-finite@cref}{{[originallemma][0][]05K0}{[1][0][]source}{}{}{}}
\makeatother
\crefname{originalsection}{original section}{original sections}
\Crefname{originalsection}{Original section}{Original sections}
\makeatletter
\newlabel{section-trace-method}{{03SH}{source}{}{sourcepointer-9}{}}
\newlabel{section-trace-method@cref}{{[originalsection][0][]03SH}{[1][0][]source}{}{}{}}
\makeatother
\crefname{originalreference}{original reference}{original references}
\Crefname{originalreference}{Original reference}{Original references}
\makeatletter
\newlabel{situation-what-to-prove}{{0A52}{source}{}{sourcepointer-10}{}}
\newlabel{situation-what-to-prove@cref}{{[originalreference][0][]0A52}{[1][0][]source}{}{}{}}
\makeatother
\title{The vanishing and surjectivity statements hold for the extension by zero of a finite locally constant sheaf of F\_ell-vector spaces along a connected open of a reduced curve}
\author{The Stacks Project Authors}
\date{stacks-project:etale-cohomology.tex:local-snapshot-20261011; GNU FDL 1.2 or later}
\begin{document}
\maketitle
\noindent\textit{Editorial note.} All delivered mathematical text is the author's own, copied from the pinned Stacks Project chapter file etale-cohomology.tex; no statement or proof step was rewritten, reordered or removed.\par\medskip
\noindent\textit{Editorial note.} The counted proof is the author's complete argument: the proof is written out in full: the trace-method reduction, the Zariski main theorem factorisation with the identification V = U x\_X Y, the shriek-pushforward compatibility, the transport along the finite morphism, and the final filtration step.\par\medskip
\noindent\textit{Editorial note.} The label lemma-vanishing-easier names the counted statement, and the proof environment immediately following it in the source is delivered verbatim as its proof.\par\medskip
\noindent\textit{Editorial note.} The counted statement is carried by the author's own lemma environment in the chapter source.\par\medskip
\noindent\textit{Editorial note.} The statement refers to the standing situation and to the enumerated statements by their labels; the eleven cited inputs are kept as external pointers with their canonical Stacks tags.\par\medskip
\noindent\textit{Allowed standard prerequisites (editorial).}\begin{itemize}
\item The trace method for a finite etale morphism of degree prime to ell. Same chapter, the section on the trace method and the lemma producing a finite etale morphism of degree prime to ell, giving maps G -> f\_*f\textasciicircum{}\{-1\}G -> G composing to an isomorphism. Supplies the reduction replacing G by f\_*f\textasciicircum{}\{-1\}G.
\item Zariski's main theorem for a quasi-finite separated morphism. More on Morphisms, the lemma factoring a quasi-finite separated morphism through a finite one by an open immersion with dense image. Supplies the diagram with the finite fbar : Y -> X and the dense open immersion j'.
\item Compatiblity of the shriek with a finite pushforward and the finite pushforward statements. Same chapter, the lemma equating the extension by zero of a finite pushforward with the finite pushforward of the extension by zero, and the lemma transporting the statements along a finite morphism. Supplies the reduction to j'\_!f\textasciicircum{}\{-1\}G and then along the finite morphism.
\item The existence of a filtration and the short exact sequence statements. Same chapter, the lemma giving a filtration of a finite locally constant sheaf by constant pieces, and the lemma on the behaviour of the statements in a short exact sequence. Supplies the final reduction to the constant sheaf Z/ellZ over a dense open of a reduced curve.
\end{itemize}\medskip
\ifdefined\hypertarget\hypertarget{sourcepointer-1}{}\fi
\noindent\textit{Source reference.} Original reference 06AP: sources/stacks-etale-cohomology/pinned-original-source.tex, source lines 15365--15365. Prerequisite pointer invoked at proof line(s) 15611; it names a standard fact of the same chapter that the counted argument uses, and every step of the counted proof built on it is retained verbatim in the delivered body.\par\medskip
\ifdefined\hypertarget\hypertarget{sourcepointer-2}{}\fi
\noindent\textit{Source reference.} Original reference 0A53: sources/stacks-etale-cohomology/pinned-original-source.tex, source lines 15335--15335. Prerequisite pointer invoked at proof line(s) 15611; it names a standard fact of the same chapter that the counted argument uses, and every step of the counted proof built on it is retained verbatim in the delivered body.\par\medskip
\ifdefined\hypertarget\hypertarget{sourcepointer-3}{}\fi
\noindent\textit{Source reference.} Original reference 0GJ1: sources/stacks-etale-cohomology/pinned-original-source.tex, source lines 11612--11612. Prerequisite pointer invoked at proof line(s) 15637; it names a standard fact of the same chapter that the counted argument uses, and every step of the counted proof built on it is retained verbatim in the delivered body.\par\medskip
\ifdefined\hypertarget\hypertarget{sourcepointer-4}{}\fi
\noindent\textit{Source reference.} Original reference 0A5D: sources/stacks-etale-cohomology/pinned-original-source.tex, source lines 15471--15471. Prerequisite pointer invoked at proof line(s) 15641; it names a standard fact of the same chapter that the counted argument uses, and every step of the counted proof built on it is retained verbatim in the delivered body.\par\medskip
\ifdefined\hypertarget\hypertarget{sourcepointer-5}{}\fi
\noindent\textit{Source reference.} Original reference 0A3R: sources/stacks-etale-cohomology/pinned-original-source.tex, source lines 10414--10414. Prerequisite pointer invoked at proof line(s) 15616; it names a standard fact of the same chapter that the counted argument uses, and every step of the counted proof built on it is retained verbatim in the delivered body.\par\medskip
\ifdefined\hypertarget\hypertarget{sourcepointer-6}{}\fi
\noindent\textit{Source reference.} Original reference 0A5C: sources/stacks-etale-cohomology/pinned-original-source.tex, source lines 15450--15450. Prerequisite pointer invoked at proof line(s) 15646; it names a standard fact of the same chapter that the counted argument uses, and every step of the counted proof built on it is retained verbatim in the delivered body.\par\medskip
\ifdefined\hypertarget\hypertarget{sourcepointer-7}{}\fi
\noindent\textit{Source reference.} Original reference 0A3Q: sources/stacks-etale-cohomology/pinned-original-source.tex, source lines 15566--15566. Prerequisite pointer invoked at proof line(s) 15649; it names a standard fact of the same chapter that the counted argument uses, and every step of the counted proof built on it is retained verbatim in the delivered body.\par\medskip
\ifdefined\hypertarget\hypertarget{sourcepointer-8}{}\fi
\noindent\textit{Source reference.} Original lemma 05K0: sources/stacks-chow/tags.txt, source lines 7184--7184. Prerequisite pointer invoked at proof line(s) 15625; it names a standard fact of a sibling chapter that the counted argument uses, and every step of the counted proof built on it is retained verbatim in the delivered body.\par\medskip
\ifdefined\hypertarget\hypertarget{sourcepointer-9}{}\fi
\noindent\textit{Source reference.} Original section 03SH: sources/stacks-etale-cohomology/pinned-original-source.tex, source lines 10330--10330. Prerequisite pointer invoked at proof line(s) 15617; it names a standard fact of the same chapter that the counted argument uses, and every step of the counted proof built on it is retained verbatim in the delivered body.\par\medskip
\ifdefined\hypertarget\hypertarget{sourcepointer-10}{}\fi
\noindent\textit{Source reference.} Original reference 0A52: sources/stacks-etale-cohomology/pinned-original-source.tex, source lines 15324--15324. Prerequisite pointer invoked at proof line(s) 15606; it names a standard fact of the same chapter that the counted argument uses, and every step of the counted proof built on it is retained verbatim in the delivered body.\par\medskip
\begin{lemma}
\label{lemma-vanishing-easier}
In Situation \ref{situation-what-to-prove} assume $X$ reduced.
Let $j : U \to X$ an open immersion with $U$ connected. Let
$\ell$ be a prime number. Let $\mathcal{G}$ a finite locally
constant sheaf of $\mathbf{F}_\ell$-vector spaces on $U$. Let
$\mathcal{F} = j_!\mathcal{G}$. Then statements
(\ref{item-vanishing}) -- (\ref{item-surjective}) hold for $\mathcal{F}$.
\end{lemma}

\begin{proof}
Let $f : V \to U$ be a finite \'etale morphism of degree prime to $\ell$
as in Lemma \ref{lemma-pullback-filtered}. The discussion in
Section \ref{section-trace-method} gives maps
$$
\mathcal{G} \to f_*f^{-1}\mathcal{G} \to \mathcal{G}
$$
whose composition is an isomorphism. Hence it suffices to prove the
lemma with $\mathcal{F} = j_!f_*f^{-1}\mathcal{G}$.
By Zariski's Main theorem
(More on Morphisms, Lemma
\ref{more-morphisms-lemma-quasi-finite-separated-pass-through-finite})
we can choose a diagram
$$
\xymatrix{
V \ar[r]_{j'} \ar[d]_f & Y \ar[d]^{\overline{f}} \\
U \ar[r]^j & X
}
$$
with $\overline{f} : Y \to X$ finite and $j'$ an open immersion
with dense image. We may replace $Y$ by its reduction (this does
not change $V$ as $V$ is reduced being \'etale over $U$).
Since $f$ is finite and $V$ dense in $Y$ we have $V = U \times_X Y$. By
Lemma \ref{lemma-compatible-shriek-push-finite} we have
$$
j_!f_*f^{-1}\mathcal{G} = \overline{f}_*j'_!f^{-1}\mathcal{G}
$$
By Lemma \ref{lemma-finite-pushforward-statements} it suffices to
consider $j'_!f^{-1}\mathcal{G}$.
The existence of the filtration given by
Lemma \ref{lemma-pullback-filtered},
the fact that $j'_!$ is exact, and
Lemma \ref{lemma-ses-statements}
reduces us to the case
$\mathcal{F} = j'_!\underline{\mathbf{Z}/\ell\mathbf{Z}}$
which is Lemma \ref{lemma-somewhat-easier}.
\end{proof}

\end{document}
